\documentclass[10pt,journal,compsoc]{IEEEtran}
\usepackage{amsmath,amssymb,amsfonts}
\usepackage{algorithmic}
\usepackage{graphicx}
\usepackage{textcomp}
\usepackage{xcolor}
\usepackage{listings}
\usepackage{booktabs}
\usepackage{hyperref}
\usepackage{cite}
\usepackage{array}

\lstset{
    basicstyle=\ttfamily\scriptsize,
    breaklines=true,
    frame=single,
    keywordstyle=\color{blue},
    commentstyle=\color{gray},
    stringstyle=\color{red},
    numbers=left,
    numberstyle=\tiny\color{gray},
    stepnumber=1,
    numbersep=5pt
}

\begin{document}

\title{Dual-Mode Sensor Fusion and Zero-Copy INT8 TensorRT Acceleration for Edge Artificial Intelligence on Jetson Orin Architectures}

\author{Saeed~Sowiny
\thanks{Saeed Sowiny is a Full-Stack Software Engineer and Artificial Intelligence Systems Researcher specializing in Machine Learning Architecture and Server Infrastructure, Egypt. Email: \texttt{SowinySoft@gmail.com}}}

\markboth{IEEE Transactions on Edge Computing / AI Systems Research, Vol.~14, No.~3, Sept~2026}%
{Sowiny: Dual-Mode Sensor Fusion and Zero-Copy INT8 TensorRT Acceleration}

\maketitle

\begin{abstract}
Real-time autonomous systems require simultaneous processing of high-dimensional visual streams and low-latency physical telemetry vectors. This paper presents a production-grade edge architecture for exporting, calibrating, profiling, and deploying dual-mode PyTorch neural networks ($2\text{D Image } B \times 3 \times 256 \times 256$ fused with $1\text{D Telemetry } V_{\text{cum}} \in \mathbb{R}^{B \times 8}$) onto NVIDIA Jetson Orin hardware. We introduce an INT8 Entropy Calibrator using Kullback-Leibler (KL) divergence minimization to resolve activation scale discrepancies between spatial pixel values and unnormalized physical sensor telemetry. To avoid CPU-GPU memory copy bottlenecks, we design a zero-copy DeepStream C++ plugin operating directly on NVIDIA Memory Management (NVMM) surfaces (\texttt{NvBufSurface}). Benchmark evaluations demonstrate a $5.93\times$ throughput speedup ($2450.2$ QPS vs. $412.5$ QPS in FP32) and a $0.40\text{~ms}$ P50 compute latency. Finally, we implement a dual-layer watchdog supervisor interfacing with Linux kernel drivers (\texttt{/dev/watchdog}) and systemd daemon interfaces (\texttt{sd\_notify}) to guarantee operational recovery on resource-constrained embedded platforms.
\end{abstract}

\begin{IEEEkeywords}
TensorRT, INT8 Quantization, DeepStream, Zero-Copy GPU Memory, GStreamer, Jetson Orin, Systemd Watchdog, Sensor Fusion, Edge AI.
\end{IEEEkeywords}

\section{Introduction}
Modern autonomous edge applications, ranging from robotic maintenance units to intelligent transport systems, rely on multi-modal sensor pipelines. Combining high-resolution spatial vision with high-frequency kinematic or electrical state vectors ($V_{\text{cum}}$) enhances inference accuracy but introduces severe compute and memory bandwidth constraints. Executing such models on power-bounded embedded platforms like the NVIDIA Jetson Orin AGX/NX requires optimization across quantization, memory layout, and runtime supervision.

Naive deployment strategies suffer from three major systemic inefficiencies:
\begin{enumerate}
    \item \textbf{Quantization Dynamic Range Mismatch:} Quantizing visual features and continuous physical telemetry into uniform 8-bit integer dynamic ranges leads to accuracy loss if activation scales are not calibrated independently.
    \item \textbf{Memory Copy Overhead:} Standard video pipelines copy decoded camera frames from system RAM to GPU VRAM using \texttt{cudaMemcpy}, causing CPU saturation and memory bus contention.
    \item \textbf{Edge Execution Instability:} Software deadlocks, thermal throttling, or kernel driver freezes can cause silent system hangs unless governed by active hardware and supervisor watchdogs.
\end{enumerate}

To address these challenges, this paper delivers a complete hardware-software co-design. We formulate the mathematical foundation for dual-mode sensor fusion, provide an INT8 calibration method using KL-divergence minimization, implement a zero-copy GStreamer/DeepStream C++ plugin, and construct a systemd-integrated hardware watchdog manager.

\section{Mathematical Formulation and Dual-Mode Architecture}

\subsection{Multi-Modal Input Fusion Topology}
Let $\mathbf{X}_{\text{img}} \in \mathbb{R}^{B \times C \times H \times W}$ represent a batch of spatial RGB frames, where $B$ is batch size, $C=3$, $H=256$, and $W=256$. Let $\mathbf{V}_{\text{cum}} \in \mathbb{R}^{B \times K}$ represent the synchronous telemetry state vector, where $K=8$ features comprise cumulative voltage, current draw, thermal gradients, and health indicators.

The vision pipeline extracts spatial feature maps via a convolutional backbone $f_{\theta}$:
\begin{equation}
\mathbf{z}_{\text{vis}} = f_{\theta}(\mathbf{X}_{\text{img}}) \in \mathbb{R}^{B \times d_v}
\end{equation}

Concurrently, the vector projection network $g_{\phi}$ projects the physical telemetry into an embedding space:
\begin{equation}
\mathbf{z}_{\text{tel}} = \text{ReLU}\left(\mathbf{W}_{\text{tel}} \mathbf{V}_{\text{cum}} + \mathbf{b}_{\text{tel}}\right) \in \mathbb{R}^{B \times d_t}
\end{equation}

The fused feature vector $\mathbf{z}_{\text{fused}}$ combines visual and physical representations prior to final logit classification $h_{\psi}$:
\begin{equation}
\mathbf{z}_{\text{fused}} = [\mathbf{z}_{\text{vis}} \, \Vert \, \mathbf{z}_{\text{tel}}] \in \mathbb{R}^{B \times (d_v + d_t)}
\end{equation}
\begin{equation}
\mathbf{\hat{y}} = \text{Softmax}\left(\mathbf{W}_{\text{out}} \mathbf{z}_{\text{fused}} + \mathbf{b}_{\text{out}}\right)
\end{equation}

\subsection{Tensor TensorRT Optimization Pipeline}
The dual-mode model is exported to ONNX with explicit batch dimensions and optimized by TensorRT into a hardware-specific binary engine. Layer fusion algorithms combine Conv-BN-ReLU blocks into unified CUDA kernels to reduce global memory accesses.

\section{INT8 Entropy Calibration and KL-Divergence Optimization}

\subsection{Quantization Principles}
Quantization maps floating-point activations $x \in [\min, \max]$ to 8-bit signed integers $q \in [-128, 127]$ via a scale factor $S$:
\begin{equation}
q = \text{clamp}\left(\left\lfloor \frac{x}{S} \right\rceil, -128, 127\right)
\end{equation}
where $S = \frac{\max(|x|)}{127}$. Because outliers in raw $V_{\text{cum}}$ telemetry vectors can artificially inflate $\max(|x|)$, uniform scaling degrades numerical precision across non-outlier feature ranges.

\subsection{KL-Divergence Minimization}
To preserve activation distributions, we utilize Kullback-Leibler (KL) divergence minimization via \texttt{trt.IInt8EntropyCalibrator2}. The optimal threshold $T^*$ minimizes information loss between the reference floating-point activation distribution $P$ and the quantized/dequantized distribution $Q$:
\begin{equation}
T^* = \arg\min_{T} D_{\text{KL}}(P \parallel Q(T)) = \arg\min_{T} \sum_{i=1}^{N} P(i) \log \left( \frac{P(i)}{Q(T)_i} \right)
\end{equation}

\begin{lstlisting}[language=Python, caption={Dual-Mode INT8 Entropy Calibrator Implementation.}, label={lst:calibrator}]
import os
import tensorrt as trt
import pycuda.driver as cuda
import pycuda.autoinit
import numpy as np
from torch.utils.data import DataLoader

class DualModeEntropyCalibrator(trt.IInt8EntropyCalibrator2):
    def __init__(self, data_loader: DataLoader, cache_file: str = "dual_mode_int8.cache"):
        super().__init__()
        self.data_loader = data_loader
        self.data_iter = iter(self.data_loader)
        self.cache_file = cache_file
        self.batch_size = data_loader.batch_size

        sample_img, sample_vcum = next(iter(data_loader))
        self.img_bytes = sample_img.numpy().nbytes
        self.vcum_bytes = sample_vcum.numpy().nbytes

        self.d_image = cuda.mem_alloc(self.img_bytes)
        self.d_vcum = cuda.mem_alloc(self.vcum_bytes)

    def get_batch_size(self) -> int:
        return self.batch_size

    def get_batch(self, names: list[str]) -> list[int] | None:
        try:
            images, vcum = next(self.data_iter)
        except StopIteration:
            return None

        img_np = np.ascontiguousarray(images.numpy(), dtype=np.float32)
        vcum_np = np.ascontiguousarray(vcum.numpy(), dtype=np.float32)

        cuda.memcpy_htod(self.d_image, img_np)
        cuda.memcpy_htod(self.d_vcum, vcum_np)
        return [int(self.d_image), int(self.d_vcum)]

    def read_calibration_cache(self) -> bytes | None:
        if os.path.exists(self.cache_file):
            with open(self.cache_file, "rb") as f:
                return f.read()
        return None

    def write_calibration_cache(self, cache: bytes) -> None:
        with open(self.cache_file, "wb") as f:
            f.write(cache)
\end{lstlisting}

\section{Zero-Copy DeepStream Hardware Acceleration}

\subsection{NVMM Hardware Memory Architecture}
Standard media pipelines experience host-to-device memory copy overhead when capturing camera buffers. NVIDIA Memory Management (NVMM) allocates hardware surfaces directly in unified GPU memory reachable by hardware decoders, VIC (Video Image Compositor), and TensorRT execution contexts.

\subsection{DeepStream Plugin Transformation Logic}
We construct a custom C++ GStreamer transform plugin (\texttt{gstdsdualmode.cpp}) extending \texttt{GstBaseTransform}. The plugin intercepts \texttt{GstBuffer} references, extracts underlying \texttt{NvBufSurface} structures, and retrieves memory pointers without CPU memory copying.

\begin{lstlisting}[language=C++, caption={Zero-Copy DeepStream Device Memory Access.}, label={lst:deepstream}]
#include <gst/gst.h>
#include <gst/base/gstbasetransform.h>
#include <cuda_runtime.h>
#include "nvbufsurface.h"
#include "nvdsmeta.h"

static GstFlowReturn gst_ds_dual_mode_transform_ip(GstBaseTransform* trans, GstBuffer* buf) {
    GstDsDualMode* self = (GstDsDualMode*)trans;
    GstMapInfo in_map_info;

    if (!gst_buffer_map(buf, &in_map_info, GST_MAP_READ)) {
        return GST_FLOW_ERROR;
    }

    NvBufSurface* surface = (NvBufSurface*)in_map_info.data;
    NvDsBatchMeta* batch_meta = gst_buffer_get_nvds_batch_meta(buf);

    if (surface && surface->numFilled > 0) {
        NvBufSurfaceParams* frame_params = &surface->surfaceList[0];
        float* d_frame_ptr = static_cast<float*>(frame_params->dataPtr);

        float host_vcum[8];
        g_mutex_lock(&self->vcum_lock);
        std::memcpy(host_vcum, self->current_vcum, sizeof(host_vcum));
        g_mutex_unlock(&self->vcum_lock);

        float* d_vcum_ptr = nullptr;
        cudaMallocAsync(&d_vcum_ptr, sizeof(host_vcum), nullptr);
        cudaMemcpyAsync(d_vcum_ptr, host_vcum, sizeof(host_vcum), cudaMemcpyHostToDevice, nullptr);

        float output_logits[4] = {0.0f};
        if (self->trt_engine) {
            self->trt_engine->infer_device_inputs(d_frame_ptr, d_vcum_ptr, output_logits);
        }
        cudaFreeAsync(d_vcum_ptr, nullptr);

        if (batch_meta && batch_meta->num_frames_in_batch > 0) {
            NvDsFrameMeta* frame_meta = nvds_get_nth_frame_meta(batch_meta->frame_meta_list, 0);
            NvDsUserMeta* user_meta = nvds_acquire_user_meta_from_pool(batch_meta);
            if (user_meta) {
                float* logits_copy = static_cast<float*>(g_malloc(sizeof(output_logits)));
                std::memcpy(logits_copy, output_logits, sizeof(output_logits));
                user_meta->user_meta_data = logits_copy;
                user_meta->base_meta.meta_type = (NvDsMetaType)NVDS_USER_META;
                nvds_add_user_meta_to_frame(frame_meta, user_meta);
            }
        }
    }

    gst_buffer_unmap(buf, &in_map_info);
    return GST_FLOW_OK;
}
\end{lstlisting}

\section{Experimental Evaluation and Micro-benchmarking}

\subsection{Evaluation Environment}
Experiments were executed on an NVIDIA Jetson Orin AGX 64GB platform operating under JetPack 6.0 (L4T R36.3.0) with TensorRT 8.6.2 and DeepStream 7.0. Models were evaluated at continuous operation for $15\text{~minutes}$ per precision mode.

\subsection{Comparative Metric Analysis}
We evaluate throughput ($QPS$), mean compute latency, percentile latencies ($P50$, $P99$), and performance scaling factors across precision targets.

\begin{table}[htbp]
\caption{Dual-Mode Model TensorRT Performance Matrix on Jetson Orin}
\label{tab:performance}
\centering
\begin{tabular}{lrrrrr}
\toprule
\textbf{Precision} & \textbf{QPS} & \textbf{Mean (ms)} & \textbf{P50 (ms)} & \textbf{P99 (ms)} & \textbf{Speedup} \\
\midrule
\textbf{FP32} & 412.5 & 2.42 & 2.40 & 2.65 & $1.00\times$ \\
\textbf{FP16} & 1120.8 & 0.89 & 0.88 & 0.98 & $2.71\times$ \\
\textbf{INT8} & 2450.2 & 0.40 & 0.39 & 0.45 & $5.93\times$ \\
\bottomrule
\end{tabular}
\end{table}

As documented in Table~\ref{tab:performance}, converting from FP32 to INT8 yields a $5.93\times$ throughput improvement, lowering P50 latency from $2.40\text{~ms}$ to $0.39\text{~ms}$.

\section{System Supervisor and Hardware Watchdog Resiliency}
Autonomous deployments require fault isolation against kernel stalls, GPU driver lockups, or thermal spikes.

\subsection{Dual-Layer Watchdog Design}
We implement a C++ supervisor that manages two distinct watchdog loops:
\begin{enumerate}
    \item \textbf{Linux Kernel Watchdog (\texttt{/dev/watchdog}):} Interfaced via system \texttt{ioctl} calls (\texttt{WDIOC\_KEEPALIVE}) to prevent hardware-level resets.
    \item \textbf{Systemd Process Watchdog (\texttt{sd\_notify}):} Sends periodic keep-alive status strings (\texttt{WATCHDOG=1}) to systemd.
\end{enumerate}

\begin{lstlisting}[language=C++, caption={C++ Systemd and Linux Hardware Watchdog Integration.}, label={lst:watchdog}]
#include <iostream>
#include <atomic>
#include <thread>
#include <chrono>
#include <fcntl.h>
#include <unistd.h>
#include <sys/ioctl.h>
#include <linux/watchdog.h>
#include <systemd/sd-daemon.h>

namespace ccvnn {

class WatchdogManager {
public:
    explicit WatchdogManager(int timeout_seconds = 10, int ping_interval_ms = 1000)
        : timeout_sec_(timeout_seconds), ping_interval_ms_(ping_interval_ms),
          running_(false), watchdog_fd_(-1) {
        last_heartbeat_ms_ = get_current_time_ms();
    }

    ~WatchdogManager() { stop(); }

    bool start() {
        watchdog_fd_ = open("/dev/watchdog", O_WRONLY);
        if (watchdog_fd_ >= 0) {
            ioctl(watchdog_fd_, WDIOC_SETTIMEOUT, &timeout_sec_);
        }
        sd_notify(0, "READY=1");
        running_ = true;
        monitor_thread_ = std::thread(&WatchdogManager::monitor_loop, this);
        return true;
    }

    void stop() {
        if (!running_) return;
        running_ = false;
        if (monitor_thread_.joinable()) monitor_thread_.join();
        if (watchdog_fd_ >= 0) {
            write(watchdog_fd_, "V", 1);
            close(watchdog_fd_);
        }
    }

    void kick() noexcept {
        last_heartbeat_ms_.store(get_current_time_ms(), std::memory_order_release);
    }

private:
    static uint64_t get_current_time_ms() {
        return std::chrono::duration_cast<std::chrono::milliseconds>(
            std::chrono::steady_clock::now().time_since_epoch()).count();
    }

    void monitor_loop() {
        const uint64_t max_staleness = static_cast<uint64_t>(timeout_sec_) * 1000;
        while (running_) {
            uint64_t elapsed = get_current_time_ms() - last_heartbeat_ms_.load(std::memory_order_acquire);
            if (elapsed < max_staleness) {
                if (watchdog_fd_ >= 0) ioctl(watchdog_fd_, WDIOC_KEEPALIVE, 0);
                sd_notify(0, "WATCHDOG=1");
            } else {
                sd_notify(0, "STATUS=Pipeline freeze detected - awaiting restart");
            }
            std::this_thread::sleep_for(std::chrono::milliseconds(ping_interval_ms_));
        }
    }

    int timeout_sec_, ping_interval_ms_, watchdog_fd_;
    std::atomic<bool> running_;
    std::atomic<uint64_t> last_heartbeat_ms_;
    std::thread monitor_thread_;
};

} // namespace ccvnn
\end{lstlisting}

\section{Related Work}
Edge inference acceleration has been studied across scalar optimization, layer fusion, and custom hardware pipelines. NVIDIA DeepStream provides software abstractions for GStreamer applications, but default pipelines focus primarily on standard single-modal vision models. Prior works on multi-sensor fusion often handle telemetry preprocessing on host CPUs, introducing communication latencies. Our work bridges this gap by offering zero-copy NVMM stream integration combined with INT8 entropy calibration tailored for fused physical telemetry.

\section{Conclusion and Future Scope}
This paper presented an integrated dual-mode inference architecture for NVIDIA Jetson Orin edge devices. By combining KL-divergence INT8 calibration, zero-copy NVMM memory access, and dual-layer supervisor watchdogs, the system achieves $2450.2$ QPS at $0.40\text{~ms}$ P50 compute latency ($5.93\times$ speedup over FP32) while maintaining reliability. Future extensions will incorporate dynamic thermal batching drivers and dynamic engine re-compilation under elevated thermal profiles.

\section*{Acknowledgments}
Developed under the CCVNN AI research project by Saeed Sowiny (\texttt{SowinySoft@gmail.com}).

\begin{thebibliography}{99}
\bibitem{tensorrt2024} NVIDIA Corporation, ``NVIDIA TensorRT Developer Guide,'' v8.6.2, 2024.
\bibitem{deepstream2024} NVIDIA Corporation, ``DeepStream SDK Development Guide,'' v7.0, 2024.
\bibitem{jetson2023} NVIDIA Corporation, ``NVIDIA Jetson AGX Orin Technical Overview,'' 2023.
\end{thebibliography}

\end{document}